\chapter{Vectors and Dot Products}
\label{ch:ch13-vectors-dot-products}

\section{From Triangles to Vector Operations}

The previous chapter concluded with an observation that deserves further
attention. The Law of Cosines,

\[
c^2=a^2+b^2-2ab\cos C,
\]

appeared at first to be a statement about triangles, yet its structure
suggests something more general. The formula relates lengths and angles in
a single expression, measuring not only how large two sides are but also
how they are oriented relative to one another. In Chapter~2, vectors were
introduced as quantities possessing both magnitude and direction, and
their lengths were computed using the distance formula. What remained
missing was an operation capable of combining two vectors while taking
both magnitude and direction into account simultaneously.

This chapter introduces such an operation. It will first appear as a
simple algebraic rule whose purpose is not immediately obvious. Only after
a derivation based on the Law of Cosines will its geometric meaning become
clear. The result is one of the most important operations in mathematics
and physics: the dot product.

\section{The Algebraic Definition}

Suppose

\[
\mathbf{u}=\langle u_1,u_2\rangle,
\qquad
\mathbf{v}=\langle v_1,v_2\rangle.
\]

\begin{definition}
The \emph{dot product} of $\mathbf{u}$ and $\mathbf{v}$ is

\[
\mathbf{u}\cdot\mathbf{v}
=
u_1v_1+u_2v_2.
\]
\end{definition}

At first glance this definition appears somewhat arbitrary. Why should the
products of corresponding coordinates be added together? Why not subtract
them, multiply them, or combine them in some other way? The definition
provides a quantity that is easy to compute, but it does not yet explain
what that quantity represents.

For example, if

\[
\mathbf{u}=\langle 3,4\rangle,
\qquad
\mathbf{v}=\langle 2,5\rangle,
\]

then

\[
\mathbf{u}\cdot\mathbf{v}
=
(3)(2)+(4)(5)
=
6+20
=
26.
\]

The computation is straightforward, but the meaning of the answer is not.
The purpose of the next section is to show that this seemingly ad hoc rule
has a remarkably natural geometric interpretation.

\section{The Geometric Formula}

The algebraic definition of the dot product provides a method of
calculation, but it does not explain why the operation is useful. To
discover its geometric meaning, consider two vectors placed with their
tails at the origin:

\[
\mathbf{u}=\langle u_1,u_2\rangle,
\qquad
\mathbf{v}=\langle v_1,v_2\rangle.
\]

Let $\theta$ denote the angle between them. The vectors and the segment
joining their endpoints form a triangle. That third side is represented by
the vector

\[
\mathbf{u}-\mathbf{v}
=
\langle u_1-v_1,\;u_2-v_2\rangle.
\]

The lengths of the three sides are therefore

\[
|\mathbf{u}|,
\qquad
|\mathbf{v}|,
\qquad
|\mathbf{u}-\mathbf{v}|.
\]

Since these lengths form the sides of a triangle, the Law of Cosines from
Chapter~12 applies directly.

\begin{theorem}[Geometric Formula for the Dot Product]
For nonzero vectors $\mathbf{u}$ and $\mathbf{v}$ with angle $\theta$
between them,

\[
\mathbf{u}\cdot\mathbf{v}
=
|\mathbf{u}|\,|\mathbf{v}|\cos\theta.
\]
\end{theorem}

\begin{proof}
Applying the Law of Cosines to the triangle formed by
$\mathbf{u}$, $\mathbf{v}$, and $\mathbf{u}-\mathbf{v}$ gives

\[
|\mathbf{u}-\mathbf{v}|^2
=
|\mathbf{u}|^2
+
|\mathbf{v}|^2
-
2|\mathbf{u}|\,|\mathbf{v}|\cos\theta.
\]

We now compute the same quantity using coordinates.

Since

\[
\mathbf{u}-\mathbf{v}
=
\langle u_1-v_1,\;u_2-v_2\rangle,
\]

the magnitude formula gives

\[
|\mathbf{u}-\mathbf{v}|^2
=
(u_1-v_1)^2+(u_2-v_2)^2.
\]

Expanding,

\[
|\mathbf{u}-\mathbf{v}|^2
=
u_1^2-2u_1v_1+v_1^2
+
u_2^2-2u_2v_2+v_2^2.
\]

Rearranging terms,

\[
|\mathbf{u}-\mathbf{v}|^2
=
(u_1^2+u_2^2)
+
(v_1^2+v_2^2)
-
2(u_1v_1+u_2v_2).
\]

Using the magnitude formula once again,

\[
u_1^2+u_2^2
=
|\mathbf{u}|^2,
\qquad
v_1^2+v_2^2
=
|\mathbf{v}|^2,
\]

so

\[
|\mathbf{u}-\mathbf{v}|^2
=
|\mathbf{u}|^2
+
|\mathbf{v}|^2
-
2(u_1v_1+u_2v_2).
\]

Both derivations describe the same quantity
$|\mathbf{u}-\mathbf{v}|^2$, so

\[
|\mathbf{u}|^2
+
|\mathbf{v}|^2
-
2(u_1v_1+u_2v_2)
=
|\mathbf{u}|^2
+
|\mathbf{v}|^2
-
2|\mathbf{u}|\,|\mathbf{v}|\cos\theta.
\]

Subtracting the common terms and dividing by $-2$ yields

\[
u_1v_1+u_2v_2
=
|\mathbf{u}|\,|\mathbf{v}|\cos\theta.
\]

Since the left-hand side is the algebraic definition of the dot product,

\[
\mathbf{u}\cdot\mathbf{v}
=
|\mathbf{u}|\,|\mathbf{v}|\cos\theta,
\]

which completes the proof.
\end{proof}

The theorem reveals why the algebraic definition was chosen. The quantity

\[
u_1v_1+u_2v_2
\]

appeared at first to be merely a computational rule, but the proof shows
that it measures something fundamentally geometric. The dot product
depends simultaneously on the lengths of the vectors and on the angle
between them. Two vectors pointing in nearly the same direction produce a
large positive dot product; vectors pointing in opposite directions
produce a negative dot product; vectors meeting at a right angle produce
zero.

The algebraic formula and the geometric formula therefore provide two
different descriptions of the same quantity. The algebraic version is
usually easier to compute from coordinates, while the geometric version
often makes interpretation easier. One of the central advantages of the
dot product is that it allows these two viewpoints to be translated into
one another.

\section{A Worked Example}

The theorem becomes more convincing when both formulas are applied to the
same pair of vectors.

\begin{example}
Let

\[
\mathbf{u}
=
\langle 3,4\rangle,
\qquad
\mathbf{v}
=
\langle 4,0\rangle.
\]

Using the algebraic definition,

\[
\mathbf{u}\cdot\mathbf{v}
=
(3)(4)+(4)(0)
=
12.
\]

Now compute the same quantity geometrically.

The magnitudes are

\[
|\mathbf{u}|
=
\sqrt{3^2+4^2}
=
5,
\]

and

\[
|\mathbf{v}|
=
\sqrt{4^2}
=
4.
\]

The angle between the vectors satisfies

\[
\cos\theta
=
\frac{3}{5},
\]

since $\mathbf{u}$ forms a $3$-$4$-$5$ triangle with the positive
$x$-axis.

The geometric formula gives

\[
\mathbf{u}\cdot\mathbf{v}
=
(5)(4)\left(\frac35\right)
=
12.
\]

Both methods produce the same result, exactly as the theorem predicts.
\end{example}

\section{Orthogonality}

One of the most important consequences of the geometric formula appears
when the angle between two vectors is a right angle.

Recall that two vectors are said to be \emph{orthogonal} if they are
perpendicular. In coordinate geometry, determining whether two vectors
meet at a right angle might appear to require finding slopes or computing
angles directly. The dot product provides a much simpler test.

\begin{theorem}
Let $\mathbf{u}$ and $\mathbf{v}$ be nonzero vectors. Then

\[
\mathbf{u}\cdot\mathbf{v}=0
\]

if and only if

\[
\mathbf{u}\perp\mathbf{v}.
\]
\end{theorem}

\begin{proof}
By the geometric formula,

\[
\mathbf{u}\cdot\mathbf{v}
=
|\mathbf{u}|\,|\mathbf{v}|\cos\theta.
\]

Since $\mathbf{u}$ and $\mathbf{v}$ are nonzero, their magnitudes are
positive, so the product can equal zero only when

\[
\cos\theta=0.
\]

This occurs precisely when

\[
\theta=90^\circ.
\]

Thus

\[
\mathbf{u}\cdot\mathbf{v}=0
\]

if and only if the vectors are perpendicular.
\end{proof}

The theorem converts a geometric question into a simple algebraic
calculation. Instead of constructing right angles or comparing slopes, one
need only compute a dot product and check whether the result is zero.

\begin{example}
Determine whether

\[
\mathbf{u}
=
\langle 2,3\rangle,
\qquad
\mathbf{v}
=
\langle 3,-2\rangle
\]

are perpendicular.

Computing the dot product,

\[
\mathbf{u}\cdot\mathbf{v}
=
(2)(3)+(3)(-2)
=
6-6
=
0.
\]

Since the dot product is zero, the vectors are orthogonal.
\end{example}

This criterion is often used in geometry, physics, computer graphics, and
engineering because it avoids measuring angles directly. Orthogonality
becomes an algebraic condition rather than a geometric construction.

\section{Finding the Angle Between Two Vectors}

The geometric formula can also be solved for the angle itself. Starting
with

\[
\mathbf{u}\cdot\mathbf{v}
=
|\mathbf{u}|\,|\mathbf{v}|\cos\theta,
\]

divide by the product of the magnitudes:

\[
\cos\theta
=
\frac{\mathbf{u}\cdot\mathbf{v}}
     {|\mathbf{u}|\,|\mathbf{v}|}.
\]

Provided neither vector is zero, the angle can then be obtained by
applying the inverse cosine function introduced in Chapter~7:

\[
\theta
=
\arccos
\left(
\frac{\mathbf{u}\cdot\mathbf{v}}
     {|\mathbf{u}|\,|\mathbf{v}|}
\right).
\]

The formula transforms a geometric measurement into a sequence of
algebraic computations.

\begin{algorithmproc}{How to Compute a Dot Product and the Angle Between Two Vectors}
Given vectors $\mathbf{u}$ and $\mathbf{v}$, first compute the dot product
using the algebraic formula
\[
\mathbf{u}\cdot\mathbf{v}
=
u_1v_1+u_2v_2.
\]
Next compute the magnitudes
\[
|\mathbf{u}|=\sqrt{u_1^2+u_2^2},
\qquad
|\mathbf{v}|=\sqrt{v_1^2+v_2^2}.
\]
Then evaluate
\[
\cos\theta
=
\frac{\mathbf{u}\cdot\mathbf{v}}
     {|\mathbf{u}|\,|\mathbf{v}|}.
\]
Finally, apply the inverse cosine function:
\[
\theta
=
\arccos
\left(
\frac{\mathbf{u}\cdot\mathbf{v}}
     {|\mathbf{u}|\,|\mathbf{v}|}
\right).
\]
\end{algorithmproc}

\begin{example}
Find the angle between

\[
\mathbf{u}
=
\langle 1,2\rangle,
\qquad
\mathbf{v}
=
\langle 3,1\rangle.
\]

First compute the dot product:

\[
\mathbf{u}\cdot\mathbf{v}
=
(1)(3)+(2)(1)
=
5.
\]

The magnitudes are

\[
|\mathbf{u}|
=
\sqrt{1^2+2^2}
=
\sqrt5,
\]

and

\[
|\mathbf{v}|
=
\sqrt{3^2+1^2}
=
\sqrt{10}.
\]

Therefore

\[
\cos\theta
=
\frac{5}{\sqrt5\,\sqrt{10}}
=
\frac{1}{\sqrt2}.
\]

Applying the inverse cosine,

\[
\theta
=
\arccos\!\left(\frac1{\sqrt2}\right)
=
45^\circ.
\]

The vectors meet at an angle of $45^\circ$.
\end{example}

\begin{practicenow}
Find the angle between

\[
\mathbf{u}
=
\langle 2,1\rangle,
\qquad
\mathbf{v}
=
\langle 1,4\rangle.
\]

Compute the dot product first, then use the geometric formula to find the
angle between the vectors.
\end{practicenow}

\section{Projection}

The geometric formula for the dot product contains more information than
the angle between two vectors. It also reveals how much of one vector
lies in the direction of another.

Suppose two vectors $\mathbf{u}$ and $\mathbf{v}$ form an angle
$\theta$. If we drop a perpendicular from the endpoint of $\mathbf{u}$ to
the line determined by $\mathbf{v}$, the resulting right triangle shows
that the length of the component of $\mathbf{u}$ in the direction of
$\mathbf{v}$ is

\[
|\mathbf{u}|\cos\theta.
\]

This quantity is called the \emph{scalar projection} of
$\mathbf{u}$ onto $\mathbf{v}$.

Because the geometric formula for the dot product is

\[
\mathbf{u}\cdot\mathbf{v}
=
|\mathbf{u}|\,|\mathbf{v}|\cos\theta,
\]

dividing by $|\mathbf{v}|$ gives

\[
|\mathbf{u}|\cos\theta
=
\frac{\mathbf{u}\cdot\mathbf{v}}
     {|\mathbf{v}|}.
\]

Thus the scalar projection can be computed directly from the dot product.

\begin{definition}
The \emph{scalar projection} of $\mathbf{u}$ onto $\mathbf{v}$ is

\[
\operatorname{comp}_{\mathbf{v}}(\mathbf{u})
=
\frac{\mathbf{u}\cdot\mathbf{v}}
     {|\mathbf{v}|}.
\]
\end{definition}

The scalar projection measures a length. Positive values indicate that
$\mathbf{u}$ points generally in the same direction as $\mathbf{v}$,
while negative values indicate that it points generally in the opposite
direction.

Often, however, a length alone is not enough. We may wish to construct
the actual vector lying along the direction of $\mathbf{v}$.

To do so, first observe that

\[
\frac{\mathbf{v}}{|\mathbf{v}|}
\]

is a unit vector pointing in the same direction as $\mathbf{v}$. If we
multiply this unit vector by the scalar projection, we obtain the desired
projection vector.

\begin{definition}
The \emph{vector projection} of $\mathbf{u}$ onto $\mathbf{v}$ is

\[
\operatorname{proj}_{\mathbf{v}}(\mathbf{u})
=
\left(
\frac{\mathbf{u}\cdot\mathbf{v}}
     {|\mathbf{v}|}
\right)
\frac{\mathbf{v}}
     {|\mathbf{v}|}
=
\frac{\mathbf{u}\cdot\mathbf{v}}
     {|\mathbf{v}|^2}
\mathbf{v}.
\]
\end{definition}

The projection vector points along $\mathbf{v}$ and has exactly the
length needed to represent the component of $\mathbf{u}$ in that
direction.

\begin{example}
Let

\[
\mathbf{u}
=
\langle 4,3\rangle,
\qquad
\mathbf{v}
=
\langle 2,0\rangle.
\]

First compute the dot product:

\[
\mathbf{u}\cdot\mathbf{v}
=
(4)(2)+(3)(0)
=
8.
\]

Since

\[
|\mathbf{v}|=2,
\]

the scalar projection is

\[
\operatorname{comp}_{\mathbf{v}}(\mathbf{u})
=
\frac{8}{2}
=
4.
\]

The vector projection is

\[
\operatorname{proj}_{\mathbf{v}}(\mathbf{u})
=
\frac{8}{2^2}
\langle 2,0\rangle
=
2\langle 2,0\rangle
=
\langle 4,0\rangle.
\]

The vector $\langle 4,0\rangle$ is precisely the horizontal component of
$\mathbf{u}$.
\end{example}

Projection provides a useful way to decompose vectors into parts. A
vector can be viewed as containing one component parallel to a given
direction and another component perpendicular to it. The projection
isolates the parallel component, leaving the remainder to account for the
perpendicular part.

\section{Applications}

The geometric interpretation of the dot product makes it useful whenever
the effect of one quantity depends on how strongly it acts in a
particular direction.

One of the most important examples occurs in physics. Suppose a force
$\mathbf{F}$ moves an object through a displacement $\mathbf{d}$. Only the
component of the force acting in the direction of motion contributes to
the work performed. A force perpendicular to the motion changes the
direction of travel but does no work.

If $\theta$ is the angle between the force and the displacement, then

\[
W
=
|\mathbf{F}|\,|\mathbf{d}|\cos\theta.
\]

Using the geometric formula for the dot product, this becomes

\[
\boxed{
W
=
\mathbf{F}\cdot\mathbf{d}.
}
\]

Work is therefore a dot product.

This interpretation explains why the cosine appears. The factor
$\cos\theta$ extracts the component of the force acting along the
direction of motion. When the force points entirely in the direction of
travel, $\theta=0^\circ$ and the work is maximized. When the force is
perpendicular to the motion, $\theta=90^\circ$ and the work is zero.

\begin{example}
A force

\[
\mathbf{F}
=
\langle 10,5\rangle
\]

moves an object through a displacement

\[
\mathbf{d}
=
\langle 4,0\rangle.
\]

The work performed is

\[
W
=
\mathbf{F}\cdot\mathbf{d}
=
(10)(4)+(5)(0)
=
40.
\]

The vertical component of the force contributes nothing because the
motion is entirely horizontal.
\end{example}

Dot products also appear in navigation, where movement in a particular
direction must be isolated from a more complicated motion; in computer
graphics, where lighting calculations depend on the angle between
surfaces and light sources; and in data analysis, where geometric ideas
about similarity are often expressed through dot products and related
operations.

In each case the underlying idea is the same. The dot product measures
not merely how large two vectors are, but how strongly they align with
one another.

\section{Looking Ahead}

This chapter completes a thread that began when vectors were introduced in
Chapter~2 and continued through the trigonometry of Chapters~8--12. The
Law of Cosines, originally derived as a statement about triangles, has
reappeared as the foundation of a general vector operation. Length,
angle, orthogonality, and projection can now all be expressed within a
single algebraic framework.

The next chapter shifts attention from vectors themselves to an
alternative coordinate system. Instead of locating points by horizontal
and vertical displacement, polar coordinates describe position using a
distance and an angle. This change of viewpoint will make rotational
structure more visible and will prepare the way for Chapters~15 and~16,
where scaling and rotation are first combined into a single operator and
then, together with vectors, absorbed into the single mathematical
language of complex numbers.

\section{Exercises}

\subsection*{Computing Dot Products}

\begin{exercise}
Compute each dot product.

\begin{enumerate}
\item[(a)]
\[
\langle 2,3\rangle\cdot\langle 4,1\rangle
\]

\item[(b)]
\[
\langle -1,5\rangle\cdot\langle 2,7\rangle
\]

\item[(c)]
\[
\langle 6,-2\rangle\cdot\langle -3,4\rangle
\]
\end{enumerate}
\end{exercise}

\begin{exercise}
Let

\[
\mathbf{u}
=
\langle 3,4\rangle,
\qquad
\mathbf{v}
=
\langle 5,2\rangle.
\]

Compute $\mathbf{u}\cdot\mathbf{v}$.
\end{exercise}

\begin{exercise}
Find the dot product of

\[
\langle 7,1\rangle
\]

and

\[
\langle -2,8\rangle.
\]
\end{exercise}

\subsection*{Geometric Interpretation}

\begin{exercise}
Use the geometric formula

\[
\mathbf{u}\cdot\mathbf{v}
=
|\mathbf{u}|\,|\mathbf{v}|\cos\theta
\]

to compute the dot product of vectors with magnitudes $5$ and $8$ that
meet at an angle of $60^\circ$.
\end{exercise}

\begin{exercise}
Vectors $\mathbf{u}$ and $\mathbf{v}$ have magnitudes

\[
|\mathbf{u}|=6,
\qquad
|\mathbf{v}|=10,
\]

and the angle between them is $120^\circ$.

Find $\mathbf{u}\cdot\mathbf{v}$.
\end{exercise}

\begin{exercise}
Verify the geometric formula for

\[
\mathbf{u}
=
\langle 3,4\rangle,
\qquad
\mathbf{v}
=
\langle 4,0\rangle,
\]

by computing the dot product both algebraically and geometrically.
\end{exercise}

\subsection*{Orthogonality}

\begin{exercise}
Determine whether the following pairs of vectors are orthogonal.

\begin{enumerate}
\item[(a)]
\[
\langle 2,3\rangle,
\qquad
\langle 3,-2\rangle
\]

\item[(b)]
\[
\langle 4,1\rangle,
\qquad
\langle 2,-8\rangle
\]

\item[(c)]
\[
\langle 5,-2\rangle,
\qquad
\langle 1,3\rangle
\]
\end{enumerate}
\end{exercise}

\begin{exercise}
Find a vector orthogonal to

\[
\langle 7,4\rangle.
\]
\end{exercise}

\begin{exercise}
Find a nonzero vector orthogonal to

\[
\langle -3,8\rangle.
\]
\end{exercise}

\subsection*{Angles Between Vectors}

\begin{exercise}
Find the angle between

\[
\langle 1,2\rangle
\]

and

\[
\langle 3,1\rangle.
\]
\end{exercise}

\begin{exercise}
Find the angle between

\[
\langle 2,1\rangle
\]

and

\[
\langle 1,4\rangle.
\]
\end{exercise}

\begin{exercise}
Find the angle between

\[
\langle 4,0\rangle
\]

and

\[
\langle 2,5\rangle.
\]
\end{exercise}

\begin{exercise}
Show that two nonzero vectors are perpendicular precisely when the angle
formula yields

\[
\theta=90^\circ.
\]
\end{exercise}

\subsection*{Projection}

\begin{exercise}
Find the scalar projection of

\[
\mathbf{u}
=
\langle 6,2\rangle
\]

onto

\[
\mathbf{v}
=
\langle 3,0\rangle.
\]
\end{exercise}

\begin{exercise}
Find both the scalar projection and vector projection of

\[
\mathbf{u}
=
\langle 4,3\rangle
\]

onto

\[
\mathbf{v}
=
\langle 2,0\rangle.
\]
\end{exercise}

\begin{exercise}
Find

\[
\operatorname{proj}_{\mathbf{v}}(\mathbf{u})
\]

for

\[
\mathbf{u}
=
\langle 5,1\rangle,
\qquad
\mathbf{v}
=
\langle 1,1\rangle.
\]
\end{exercise}

\begin{exercise}
A vector is decomposed into a component parallel to a given direction and
a component perpendicular to it. Explain how projection provides the
parallel component.
\end{exercise}

\subsection*{Applications}

\begin{exercise}
A force

\[
\mathbf{F}
=
\langle 8,3\rangle
\]

moves an object through a displacement

\[
\mathbf{d}
=
\langle 5,0\rangle.
\]

Find the work performed.
\end{exercise}

\begin{exercise}
A force

\[
\mathbf{F}
=
\langle 12,5\rangle
\]

acts through a displacement

\[
\mathbf{d}
=
\langle 3,4\rangle.
\]

Compute the work.
\end{exercise}

\begin{exercise}
Explain why a force perpendicular to a displacement performs no work.
\end{exercise}

\begin{exercise}
A boat moves due east while a current pushes due north. Explain, using
the dot product, why the current contributes no work in the direction of
travel.
\end{exercise}

\subsection*{Mixed Problems}

\begin{exercise}
For

\[
\mathbf{u}
=
\langle 2,5\rangle,
\qquad
\mathbf{v}
=
\langle 4,-1\rangle,
\]

compute:

\begin{enumerate}
\item[(a)] the dot product,
\item[(b)] the angle between the vectors,
\item[(c)] the scalar projection of $\mathbf{u}$ onto $\mathbf{v}$.
\end{enumerate}
\end{exercise}

\begin{exercise}
Determine whether the triangle formed by the vectors

\[
\langle 4,1\rangle
\]

and

\[
\langle 1,-4\rangle
\]

contains a right angle.
\end{exercise}

\begin{exercise}
Show that

\[
\mathbf{u}\cdot\mathbf{u}
=
|\mathbf{u}|^2
\]

for every vector $\mathbf{u}$.
\end{exercise}

\begin{exercise}
Use the result of the previous exercise to derive the magnitude formula

\[
|\mathbf{u}|
=
\sqrt{\mathbf{u}\cdot\mathbf{u}}.
\]
\end{exercise}

\subsection*{Conceptual and Exploratory Problems}

\begin{exercise}
The algebraic definition of the dot product appears to depend on
coordinates, while the geometric formula refers only to lengths and
angles. Explain why the theorem relating the two descriptions is
important.
\end{exercise}

\begin{exercise}
Why does the dot product become negative when the angle between two
vectors exceeds $90^\circ$?
\end{exercise}

\begin{exercise}
Describe the geometric meaning of the quantity

\[
|\mathbf{u}|\cos\theta.
\]
\end{exercise}

\begin{exercise}
Suppose two vectors have fixed lengths. How does their dot product change
as the angle between them increases from $0^\circ$ to $180^\circ$?
\end{exercise}

\begin{exercise}
The geometric formula for the dot product was derived from the Law of
Cosines. Explain how this chapter extends the results of Chapter~12 from
triangles to vectors.
\end{exercise}

\subsection*{Challenge Problems}

\begin{exercise}
Prove the geometric formula for the dot product by beginning with

\[
|\mathbf{u}+\mathbf{v}|^2
\]

instead of

\[
|\mathbf{u}-\mathbf{v}|^2.
\]
\end{exercise}

\begin{exercise}
Show that

\[
|\mathbf{u}-\mathbf{v}|^2
=
|\mathbf{u}|^2
+
|\mathbf{v}|^2
-
2(\mathbf{u}\cdot\mathbf{v}).
\]

Explain how this identity reflects the proof of the geometric formula.
\end{exercise}

\begin{exercise}
Find all vectors

\[
\langle x,y\rangle
\]

that are orthogonal to

\[
\langle 3,4\rangle.
\]
\end{exercise}

\begin{exercise}
Suppose the algebraic definition of the dot product did not agree with
the geometric formula. What geometric properties of vectors would fail?
Discuss how orthogonality, angle measurement, and projection would be
affected.
\end{exercise}
